\section{Safety and Diagnostics}\label{sec:security}

A \LaTeX{} file is a program, not just a document. Compiling one with a
traditional \TeX{} engine can read arbitrary files, write new ones and, with
shell escape enabled, run commands on your machine. \texdocx{} is designed to
convert documents you did not write: papers from co-authors, submissions,
templates downloaded from the web. It therefore treats every input as
untrusted. It does not call a \TeX{} engine at all; it reimplements the
\LaTeX{} front end (tokenizer, macro expansion, document structure) and only
ever produces a \texttt{.docx} file.

\subsection{Sandbox guarantees}\label{sec:security-sandbox}

The following properties hold for every document, whatever macros it defines.
The project's security test suite checks the main ones on every test run.

\begin{itemize}
  \item \textbf{Root confinement.} \cmd{input}, \cmd{include}, \cmd{subfile}
    and \cmd{includegraphics} can read only files inside the project folder:
    by default the folder of the main file, or the folder given with
    \verb|--root|. Absolute paths, paths starting with \verb|~| and \verb|../|
    escapes are refused with error E004. Such paths are rejected on their text
    alone, before any file is looked up, so the file system outside the
    project is not touched. The
    refusal reads the same whether the outside file exists or not, so a
    document cannot use it to discover your files.
  \item \textbf{Symbolic links.} Every candidate file is resolved to its real
    path before it is opened. A link inside the project that points outside
    it is refused exactly like a \verb|../| path.
  \item \textbf{Restricted file types.} Source commands read
    \texttt{.tex}, \texttt{.ltx}, \texttt{.bbl}, \texttt{.tikz} and
    \texttt{.pgf} files; \cmd{includegraphics} reads only image formats
    (PNG, JPEG, GIF, BMP, SVG, PDF, EPS, PS). A file with any other extension,
    such as \texttt{.key} or \texttt{.json}, is refused with E004 even when
    it lies inside the project. A file with \emph{no} extension is still
    read, as in \LaTeX{}, where \verb|\input{chapter}| falls back to a file
    named \texttt{chapter}; at present this also lets dotfiles such as
    \texttt{.env} through.
  \item \textbf{No shell.} \cmd{write18} is recognised and ignored with
    warning W017; no command is ever executed. For the same reason, PDF and
    EPS figures are not converted with external tools: \texdocx{} uses a PNG,
    SVG or JPEG file of the same name if one exists, and otherwise inserts a
    placeholder.
  \item \textbf{No file writes.} \cmd{openout} is ignored with warning W016.
    The only file written is the \texttt{.docx} you asked for.
  \item \textbf{No network.} A file name that looks like a URL is refused as a
    remote resource. Links made with \cmd{href} and \cmd{url} are copied
    into the Word file as hyperlinks; they are never fetched.
  \item \textbf{Resource budgets.} Macro expansion is limited to 5\,000\,000
    steps, a nesting depth of 10\,000, 50\,000\,000 tokens and 60 seconds of
    wall-clock time. Each file the document reads may be at most 20\,MiB,
    and one run reads at most 2\,000 files. A macro bomb such as \verb|\def\a{\a\a}\a| stops with
    error E003, which names the macro stack that ran away; the text converted
    up to that point is kept.
\end{itemize}

Text from the document is always escaped before it reaches the Word XML, so
words such as \texttt{INCLUDETEXT} or \texttt{DDE} in the source remain plain
text and cannot turn into active Word fields.

\begin{tip}
The sandbox protects everything outside the project folder. Treat the folder
itself as readable by the document, because the type check does not stop
files without an extension: keep credentials, \texttt{.env} files and
private keys somewhere else, and point \verb|--root| at the smallest folder
that holds the sources and figures.
\end{tip}

\subsection{Why there is no shell at all}\label{sec:security-shell}

Many \TeX{} installations allow a restricted shell escape that runs only a
list of trusted programs. That design depends on every engine enforcing the
list correctly, and it has failed in practice.
\href{https://nvd.nist.gov/vuln/detail/cve-2023-32700}{CVE-2023-32700}
describes how LuaTeX before version 1.17.0 let a document reach the original
\texttt{io.popen} function and so run arbitrary shell commands while
compiling an untrusted file. The National Vulnerability Database rates it
CVSS~3.1 base score 7.8 (High). The fix is to upgrade to LuaTeX 1.17.0 or
later; for distributions, that means TeX Live 2023 from revision 66984 on,
or MiK\TeX{} 23.5 and newer. \texdocx{} avoids this whole class of
problem by having no code path that starts a process: there is no allow-list
to bypass.

\subsection{Risks and recommended actions}\label{sec:security-risks}

Table~\ref{tab:security-risks} lists the main risks of converting an
untrusted document and what to do about each. The severity column is the
\texdocx{} project's own rating of the impact if the risk were not handled;
only the CVE row carries an official CVSS score.

\begin{table}[ht]
  \centering
  \caption{Risks of converting untrusted \LaTeX{} and how they are handled.}
  \label{tab:security-risks}
  \begin{tabularx}{\textwidth}{XlX}
    \toprule
    \textbf{Risk} & \textbf{Severity} & \textbf{Recommended action} \\
    \midrule
    Shell escape (\cmd{write18}, CVE-2023-32700) & High & None in \texdocx{} (never run, W017); patch \TeX{} installs \\
    Reading files outside the project & High & Keep \verb|--root| at the project folder; E004 blocks escapes \\
    Reading secrets inside the project (dotfiles, files without extension) & Medium & Keep secrets out of the \verb|--root| folder \\
    Symbolic links leaving the project & High & None needed; real paths are checked \\
    Writing files (\cmd{openout}) & Medium & None needed; ignored with W016 \\
    Fetching remote resources & Medium & None needed; URLs are refused with E004 \\
    Macro bombs, runaway recursion & Medium & Rely on budgets (E003); add a process timeout on servers \\
    Silently lost content & Low & Review W013 and W014, or use \verb|--strict| \\
    \bottomrule
  \end{tabularx}
\end{table}

\subsection{Reading diagnostics}\label{sec:security-diagnostics}

\texdocx{} reports problems in the style of a modern compiler: a severity and
a stable code, a message, the file, line and column, the offending source
line with a caret under the exact spot, and often a hint. The following
output comes from a short file that uses an undefined macro and tries to
read a file outside its folder:

\begin{verbatim}
warning[W013]: unknown command \highlightbox, rendered its arguments as text
 --> report.tex:3:13
  |
3 | Results are \highlightbox{significant} at $p<0.05$.
  |             ^^^^^^^^^^^^^
  = hint: define it with \newcommand or add a package handler

error[E004]: \input{../secret}: the file is outside the project folder
 --> report.tex:4:1
  |
4 | \input{../secret}
  | ^
\end{verbatim}

The run still writes a \texttt{.docx}, but exits with status 1 because there
was an error. With \verb|--json| each diagnostic is printed as one JSON
object per line, which is convenient for editors and continuous integration;
\verb|--strict| turns every warning into an error. Three codes are worth
knowing:

\begin{description}
  \item[W013 (unknown command)] \texdocx{} has no handler for the command and
    no definition was found. Its brace arguments are kept as ordinary text,
    but an optional argument in square brackets is dropped, and the
    formatting the command would have applied is lost. Define the command with \cmd{newcommand} in your preamble, or check
    for a typo.
  \item[W014 (unsupported feature approximated)] The construct was understood
    but has no exact Word equivalent, so \texdocx{} approximated it. Examples
    are a PDF figure replaced by a placeholder, \cmd{hfill} turned into a tab,
    or an image rotation that is not applied. Read the message and check that
    spot in the Word file.
  \item[E004 (file access denied)] A file request broke a sandbox rule: it
    was outside the project, an absolute path, a URL, a file type the
    command may not read, or over a size or file-count limit. The file is never opened: a refused \cmd{input} inserts nothing, and a
    refused figure appears as a placeholder that names the file.
    Move the file into the project folder or change \verb|--root|.
\end{description}
